% =============================================================================
% PAPER CAVEAT: baseline pre-training budget
%
% For the PUBLISHED paper. The follow-up work (converged budgets, heterogeneous
% observations, the HyperBO-kernel hybrid) is deliberately NOT referenced here --
% it belongs in the follow-up paper.
%
% All numbers below are measured, not estimated. Matched pairs from the same runs
% (n = 68 at 2k iters, n = 14 at 25k), medians:
%
%     budget            EM pre-train   HyperBO pre-train   ratio
%     2,000 iters           24.6 s          405.0 s        16.4x
%     25,000 iters          20.1 s         3471.6 s       173.1x
%
% Two things this caveat must do, and one it must not:
%   DO  state the budget plainly, so the comparison is reproducible;
%   DO  note that the budget already costs the baseline 16x our pre-training time,
%       which is why it was chosen and why raising it is not free;
%   DON'T claim the ordering is unchanged at larger budgets -- we now know it is not,
%       and asserting otherwise in the published paper would be wrong.
% =============================================================================

% ---------------------------------------------------------------------------
% OPTION A -- for the experimental-setup section (recommended).
% ---------------------------------------------------------------------------

The meta-learning baselines (HyperBO, PACOH) are pre-trained for $2{,}000$ optimizer
steps. At this budget HyperBO's pre-training already costs roughly $16\times$ the
Empirical GP's ($405$s versus $25$s median on our hardware), since the EM pre-training
is a closed-form fixed-point iteration rather than a gradient-based fit. We note that
deep-kernel meta-learners are sensitive to this budget in ways that are not visible in
converged metrics, and that a definitive comparison would tune the pre-training budget
of every method to convergence rather than fixing it in advance; we leave that to future
work. The results reported here should therefore be read as a comparison at a fixed,
equal-step budget, not as a characterization of each method's best attainable
performance.

% ---------------------------------------------------------------------------
% OPTION B -- compressed, for a limitations paragraph.
% ---------------------------------------------------------------------------

Baselines are pre-trained for a fixed $2{,}000$ steps, which already costs HyperBO
$\sim\!16\times$ the Empirical GP's pre-training time. Deep-kernel meta-learners are
budget-sensitive, so these results characterize a fixed-budget comparison rather than
each method's best attainable performance; tuning every method's pre-training budget to
convergence is left to future work.

% ---------------------------------------------------------------------------
% OPTION C -- one sentence, if space is critical.
% ---------------------------------------------------------------------------

Baselines use a fixed $2{,}000$-step pre-training budget -- already $\sim\!16\times$ the
Empirical GP's pre-training cost -- and a definitive comparison would tune each method's
budget to convergence.

% ---------------------------------------------------------------------------
% WHAT NOT TO WRITE
% ---------------------------------------------------------------------------
% Avoid any variant of "we verified the ordering is robust to the pre-training
% budget." Internally it is not: at a converged 25k-step budget HyperBO's 7D
% regression NLL improves substantially and the ordering changes. The caveat above is
% honest precisely because it declines to make that claim -- it says the comparison is
% at a fixed budget and stops there, which is both true and defensible if a reader
% later runs the larger budget themselves.
